\section{Konteks Kurikulum OBE}

\subsection{Apa itu Outcome-Based Education?}

Outcome-Based Education (OBE) didefinisikan oleh Spady sebagai pendekatan yang mengorganisasi segala sesuatu dalam sistem pendidikan di sekitar hal-hal yang esensial bagi semua siswa untuk mampu melakukannya dengan sukses pada akhir pengalaman belajar \cite{spady1994}. Fokus bergeser dari input (bahan ajar, jam tatap muka) menuju output (kompetensi yang dapat didemonstrasikan). Prinsip ini menuntut perumusan learning outcomes yang jelas, desain asesmen yang valid, serta strategi pembelajaran yang mendukung pencapaian kompetensi tersebut.

\subsection{Komponen Kerangka OBE}

Kerangka OBE mencakup empat komponen dasar yang saling terhubung \cite{obe_springer2024}. Pertama, learning outcomes harus dirumuskan secara spesifik dan terukur. Kedua, sumber belajar dan materi ajar diidentifikasi untuk mendukung pencapaian outcomes. Ketiga, metode pembelajaran dipilih agar sesuai dengan tipe kompetensi yang ditargetkan. Keempat, asesmen dirancang untuk mengevaluasi tingkat pencapaian outcomes melalui penilaian formatif maupun sumatif. Keempat komponen ini membentuk alignment yang menjadi inti implementasi OBE.

\subsection{Relevansi dengan Kurikulum Program Studi}

Buku ajar Logika Matematika ini mendukung kurikulum Program Studi Teknik Informatika dengan berkontribusi terhadap Capaian Pembelajaran Lulusan (CPL) dalam aspek pengetahuan, keterampilan, dan sikap. Secara pengetahuan, mahasiswa memahami struktur logika formal, proposisi, operator, predikat, dan metode pembuktian. Secara keterampilan, mahasiswa mampu menerjemahkan pernyataan ke bentuk logika, membangun bukti, dan menulis program Prolog sederhana. Secara sikap, pendekatan sistematis dalam logika mendorong ketelitian dan konsistensi dalam berpikir.

Pemetaan eksplisit materi ke CPMK dan Sub-CPMK memastikan bahwa setiap topik dalam buku terhubung langsung dengan indikator kompetensi. Asesmen---meliputi latihan, kuis, proyek, dan ujian---dirancang untuk mengukur pencapaian indikator tersebut, sehingga dosen dan mahasiswa dapat mengevaluasi kemajuan pembelajaran secara obyektif.

Kurikulum OBE yang diterapkan Program Studi relevan dengan standar industri dan kebutuhan lulusan Teknik Informatika. Kemampuan berpikir logis, merancang argumen valid, dan memahami dasar-dasar reasoning formal menjadi fondasi bagi pengembangan perangkat lunak, kecerdasan buatan, basis data, dan sistem kritis lainnya \cite{benari_matlogic}. Dengan demikian, mata kuliah Logika Matematika tidak sekadar memenuhi requirement kurikulum, melainkan menyiapkan mahasiswa menghadapi tuntutan dunia kerja dan studi lanjut.
